\ifJA

\section{こう考えて読んでみてください}
\label{dan:000}

『伊勢物語』は、ときどき「ことば足らず」な作品だと言われます。
しかし、本当にことばが足りないのでしょうか。

人物がなぜそうしたのか。
誰がそこにいたのか。
なぜ次の行動が起こったのか。
『伊勢物語』は、そうしたことをいちいち説明してくれません。

けれども、少し考え方を変えてみてください。

\begin{quote}
\textbf{いったん出てきた人や物を、
書かれていないからといって消さない。}
\end{quote}

誰がどこにいるのか。
そこには何があるのか。
何時ごろなのか。
誰に見られているのか。
その人物には何ができて、何ができないのか。

それまでに置かれたものを、その場に残したまま
次の文を読んでみてください。

すると、「書いてないから分からない」と思っていたところが、
案外、書かなくても分かる程度になっていることがあります。

そして、もう一つ。

\begin{quote}
\textbf{書いていない理由を、急いで補わないでください。}
\end{quote}

「悲しかったから泣いた」
「恥ずかしかったから帰った」
「心が通じていなかったから話せなかった」。

そう説明すれば、話はすぐ分かりやすくなります。
でも、本文がそう言っていないのなら、
まずは説明しないまま読んでみます。

人物が何をしたかだけを追っていると、
「この人、かわいいな」
「この男、なかなかやるな」
「なんで今そこで来るんだよ」と、
こちらの側に人物像ができてくることがあります。

それでいいのです。

この本の題名を『読み手勝手伊勢物語』としたのは、そのためです。

\begin{quote}
\textbf{訳者が勝手に決めないから、読み手が勝手に読める。}
\end{quote}

ここにある訳も、できるだけ一つの読みを
正解として埋め込まないようにしました。

分からないところは、分からないまま残します。
複数に読めるところは、なるべく複数に読めるまま残します。

まず、人と物をその場に置いてみてください。

そして、その場で何ができるのかを考えてみてください。

『伊勢物語』は、思っていたほど
「ことば足らず」ではないかもしれません。

\else

\section{Try Reading It This Way}
\label{dan:000}

\textit{The Tale of Ise} is sometimes said to leave too much unsaid.
But does it really?

Why did someone do what they did?
Who was there?
Why did the next action occur?
\textit{The Tale of Ise} does not explain such things one by one.

Try approaching it a little differently.

\begin{quote}
\textbf{Once a person or an object has appeared,
do not make it disappear simply because it is no longer mentioned.}
\end{quote}

Who is where?
What is there?
What time is it?
Who can see whom?
What can a person do there, and what can they not do?

Keep what has already been placed in the scene where it is,
and then read the next sentence.

Something that first seemed impossible to understand
because ``the text doesn't say''
may turn out not to require saying at all.

And one more thing.

\begin{quote}
\textbf{Do not rush to supply a reason the text does not give.}
\end{quote}

``She cried because she was sad.''
``He left because he was embarrassed.''
``They could not speak because they did not understand each other.''

Explanations like these quickly make a story easier to understand.
But if the text does not say so,
try reading without supplying the explanation first.

Follow only what the characters actually do,
and you may find yourself thinking,
``I like this person,''
``This man is pretty good,''
or
``Why show up now?''

A character begins to take shape on your side of the text.

That is enough.

This is why the Japanese title of this book is
\textit{Yomite Katte Ise Monogatari}.

\begin{quote}
\textbf{The less the translator decides,
the more the reader is free to read.}
\end{quote}

The translations in this book likewise try not to build
one particular reading into the text as the correct answer.

What remains unclear may remain unclear.
Where more than one reading is possible,
we try to leave more than one reading possible.

First, put the people and things in place.

Then ask what can happen there.

\textit{The Tale of Ise} may not leave
as much unsaid as we thought.

\fi
